\section{Introduction}\label{sec:introduction}

Public protection against production disruptions can have a countervailing effect on industrial resilience. By making a firm's primary plant more reliable, it reduces the private return to maintaining backup production elsewhere. When jurisdictions also compete to attract mobile plants, each government may value the resulting relocation of primary production even if the induced decline in contingency readiness makes additional protection unattractive from a national perspective. The policy question is therefore not only whether local governments provide too much protection, but also how the distortion changes when protection becomes more valuable and firms can adjust both their backup capacity and their locations.

We study this question in a sequential game with two jurisdictions and two differentiated-product firms. Governments first choose risk-reducing local public inputs. Firms then select their primary locations using privately observed location-fit shifters, choose costly contingency readiness, and compete \`a la Cournot after localized production disruptions. Co-location generates correlated primary-plant failures, and product-market interaction affects the value of surviving when a rival cannot produce. Throughout, the coordinated benchmark selects only the public protection levels: firms continue to choose locations, readiness, and product-market quantities privately. The national objective measures industrial surplus net of resilience and public-investment costs. It is not a first-best allocation of all private decisions or a comprehensive valuation of disaster protection.

The first result establishes a full-game conflict between the industrial value of protection and local policy incentives. Public protection raises primary reliability directly and saves some privately incurred readiness costs, but it also crowds out backup capacity whose social availability value firms do not fully internalize. These opposing effects can make the national industrial marginal value of protection negative. A jurisdiction can nevertheless gain from unilateral protection because a safer primary site attracts mobile production and location-attached benefits. Under an exact rational benchmark and throughout a nonempty open set of the maintained symmetric primitive family, the coordinated planner uniquely chooses zero additional protection, while the decentralized game has a strictly positive symmetric Nash equilibrium. The statement concerns global policy choices, not merely signs of first-order conditions: the private continuation is uniquely determined on the policy domain, and the local equilibrium policies withstand feasible unilateral deviations. The marginal decomposition $M_L=M_S/4+2b\lambda_0$ isolates the attraction incentive under the baseline welfare-incidence assumption; the identity itself is not claimed as a new fiscal-competition theorem.

A zero coordinated optimum, however, could make the result appear to be a special corner solution. We therefore introduce a commonly valued direct benefit of protection, including nonindustrial value omitted from the initial industrial-surplus benchmark. This extension admits a strictly positive, unique coordinated protection optimum while decentralized governments still invest more. More importantly, it identifies a comparative-static question that is not answered simply by the existence of fiscal overprovision: \emph{does an increase in the common direct benefit widen or narrow the gap between decentralized and coordinated protection?} The gap need not move in the same direction as either policy level.

Our central additional result isolates the interaction between two private adjustment margins. In a separable-product benchmark with both contingency readiness and primary location endogenous, a larger direct benefit strictly \emph{widens} the decentralized-minus-coordinated policy gap. We then re-solve the complete first-stage policy game under two explicit restrictions, maintaining each restriction after every possible government deviation. If readiness is fixed while locations are reoptimized, the benefit-induced gap response is \emph{zero}. If primary locations are fixed while readiness is reoptimized, the response is \emph{negative}. These positive--zero--negative signs follow from exact, globally solved policy comparisons at the specified primitives; they are not inferred by deleting terms from an equilibrium first-order condition. They identify a within-model interaction in the \emph{comparative static of the policy gap}, not a general claim that both private decisions are necessary for overinvestment to exist.

The analytical results extend beyond a single numerical configuration. Strict sufficient inequalities in the risk, readiness-cost, location-dispersion, hosting-benefit, and protection-cost primitives guarantee the relevant global policy comparisons in the separable-product economy. In particular, where the location-attached benefit exceeds location-fit dispersion under the other maintained conditions, the overprovision gap increases with the common direct protection benefit. The positive coordinated optimum, the decentralized ranking, and this positive gap response persist on a nonempty open neighborhood with strictly positive product substitutability. A separate exact right-derivative calculation shows that greater product substitutability can itself increase the policy gap near the separable limit, giving the Cournot component a policy-level comparative static rather than merely a richer continuation game. Small changes in regional industrial-surplus incidence and in whether local hosting rewards depend on actual primary-plant operation preserve the ranking locally. None of these statements asserts a universal sign for every admissible primitive vector.

These contributions build on well-established public-economics mechanisms. Competition for mobile activity can distort productive public inputs \citep{KeenMarchand1997,BayindirUpmann1998}; firm mobility, oligopolistic location, and market power change government incentives \citep{WalzWellisch1996,MaurerWalz2000,Matsumoto1998,Billon2023,MoritaOkoshi2025}. Private geographic redundancy also entails a wedge between private and social incentives \citep{GrossmanHelpmanLhuillier2023}. Thus neither public-input overprovision, plant attraction, nor an inefficient private resilience choice is independently novel. The contribution is instead a set of globally justified policy comparisons that connects these mechanisms and establishes how their joint operation changes the \emph{response} of decentralized overprovision to a common protection benefit.

Research on disaster protection has already emphasized endogenous exposure and private adjustment. Public protection can attract investment to risky locations \citep{KouskyLuttmerZeckhauser2006}, anticipated compensation can distort prevention and residential location \citep{SeshimoYamazaki2018}, and public storm protection can affect private self-protection and insurance \citep{MahmudBarbier2016}. The wider safe-development literature similarly considers protection-induced exposure and preparedness \citep{HaerEtAl2020,FusinatoEtAl2024}. These studies caution against attributing protection's welfare effects solely to its engineering performance. Our narrower question is how competing jurisdictions choose primary-site risk reduction after firms' geographically separate production-continuity investments and primary locations respond, and how the resulting policy gap changes across fully specified counterfactual games. Section~\ref{sec:literature} develops the closest comparisons and the limits of the novelty claim.

The results have an implication for evaluating public investments in industrial resilience: changes in primary-plant reliability should be assessed alongside the endogenous substitution away from backup production and the transfers of mobile production between jurisdictions. Because the coordinated benchmark retains decentralized firm decisions, it does not automatically eliminate the underlying private--social readiness wedge. Nor does the benchmark cover the value of avoided losses to life, housing, infrastructure, or unrelated sectors. The theory establishes possibilities and conditional comparative statics, not the frequency or optimal scale of real-world disaster-protection projects.

Section~\ref{sec:model} introduces the primitives, and Section~\ref{sec:equilibrium} solves the product-market, backup, and location continuations. Section~\ref{sec:main} establishes the baseline policy conflict, and Section~\ref{sec:welfare} explains the welfare accounting and readiness wedge. Section~\ref{sec:benchmarks} develops the mechanism benchmarks, while Section~\ref{sec:analytic-region} gives primitive sufficient conditions, the response-interaction comparisons, and the differentiated-product extensions. Sections~\ref{sec:institutions} and \ref{sec:literature} discuss interpretation and related research; Sections~\ref{sec:discussion} and \ref{sec:conclusion} address scope and implications. Appendix~\ref{app:verification} contains the mathematical proofs and computational certificate descriptions.
